%%%%%%%%%%%%%%%%%%%%%%%%%%%%%%%%%%%%%%%%%
% Medium Length Professional CV
% LaTeX Template
% Version 2.0 (8/5/13)
%
% This template has been downloaded from:
% http://www.LaTeXTemplates.com
%
% Original author:
% Trey Hunner (http://www.treyhunner.com/)
%
% Important note:
% This template requires the resume.cls file to be in the same directory as the
% .tex file. The resume.cls file provides the resume style used for structuring the
% document.
%
%%%%%%%%%%%%%%%%%%%%%%%%%%%%%%%%%%%%%%%%%

%----------------------------------------------------------------------------------------
%	PACKAGES AND OTHER DOCUMENT CONFIGURATIONS
%----------------------------------------------------------------------------------------

\documentclass{resume} % Use the custom resume.cls style
\usepackage{marvosym}
\usepackage{fontawesome}
\usepackage[hidelinks,colorlinks=true,linkcolor=blue,citecolor=blue,urlcolor  = blue]{hyperref}
\usepackage[left=0.4 in,top=0.2in,right=0.4 in,bottom=0.3in]{geometry} % Document margins
\newcommand{\tab}[1]{\hspace{.1\textwidth}\rlap{#1}} 
\newcommand{\itab}[1]{\hspace{0em}\rlap{#1}}
\name{Fabrizio Solf Rivas} % Your name
%\address{Calle los Cocos lt. 44, Cieneguilla \\ \Telefon\ (+51) 980621944} % Your address
%\address{123 Pleasant Lane \\ City, State 12345} % Your secondary addess (optional)
\address{\faLinkedin\ \href{https://www.linkedin.com/in/fsolfsr/}{linkedin.com/in/fsolfsr/} \\ \Letter\ \href{mailto:fabrizio.solfrivas@studbocconi.it}{fabrizio.solfrivas@studbocconi.it} \\ \faGithub\ \href{https://github.com/fabrizio-solf}{github.com/fabrizio-solf} \\ \Telefon\ (+39) 3458820375}  % Your phone number and email
\begin{document}

%----------------------------------------------------------------------------------------
%	OBJECTIVE
%----------------------------------------------------------------------------------------
\begin{rSection}{PROFILE}

{I am an MSc student in Economic and Social Sciences at Bocconi with an interest in applying mathematical and statistical methods to questions in macroeconomics and finance. My background combines research experience in economics with data science experience in finance.}

%{I am currently pursuing a Master’s in Economics at Bocconi University (Italy) and hold a degree in Economics from the Universidad del Pacfico (Peru). I have three years of combined experience as a data scientist and research assistant. My motivation for theory and programming helped me win the award for the best bachelor's thesis in macroeconomics granted by the Central Bank of Peru. I also completed the MITx MicroMasters in Statistics. My main motivation is to develop quantitative models and apply advanced mathematical methods to address economic problems.}

\end{rSection}

%----------------------------------------------------------------------------------------
%	EDUCATION SECTION
%----------------------------------------------------------------------------------------

\begin{rSection}{Education}

{\bf MSc in Economics and Social Sciences} \hfill {Sep 2025 - Jun 2027}
\\
Bocconi University
%Master focused on quantitative methods and economic theory.

{\bf MicroMasters program in Statistics and Data Science} \hfill {Sep 2023 - Apr 2025}
\\
Massachusetts Institute of Technology (MIT)
%Courses: Data Analytics, Probability, Machine Learning, Statistics, and Statistical Modeling.

{\bf Bachelor's in Economics} \hfill {Mar 2018 - Dec 2022}
\\
Universidad del Pacífico 
%\\
%Top 5\% of the program: 2022.


\end{rSection}


%----------------------------------------------------------------------------------------
%	WORK EXPERIENCE
%----------------------------------------------------------------------------------------

\begin{rSection}{Research and Work Experience}
{\bf Bocconi University - Research Assistant} \hfill{Sep 2026 - Nov 2026}
\begin{itemize}
\setlength\itemsep{-0.5em}
    \item Assisted in research project \textbf{The Economic Costs of Ambiguous Laws} \href{https://www.ifo.de/DocDL/cesifo1_wp11929.pdf}{(Giommoni et al., 2025)}. Work with geographical information and obtained 1.3 million coordinates for firms. Used discontinuity regression analysis.  
\end{itemize}

{\bf Innocenzo Gasparini Institute for Economic Research - Visiting Student} \hfill {Apr 2026 - Present}
\begin{itemize}
\setlength\itemsep{-0.5em}
    \item  Visiting Student in a competitive research training program involving seminars and faculty mentorship
\end{itemize}


{\bf Banco de Crédito del Perú (BCP) - Senior Risk Data Scientist} \hfill {Feb 2025 - Aug 2025}
\begin{itemize}
\setlength\itemsep{-0.5em}
    %\item Credit Risk Modeler at the biggest retail bank in Peru.
    \item Developed machine learning models of household financial vulnerability and over-indebtedness using large-scale financial data, with emphasis on prediction, validation, and model interpretability. %Algorithms: {\it gradient boosting}.
\end{itemize}


{\bf Ernst \& Young (EY) - In-Charge Data Scientist} \hfill {May 2023 - Jan 2025}
\begin{itemize}
\setlength\itemsep{-0.5em}
    %\item Leader of the Analytics Lab of the Pricing COE.
    %\item I lead the project to design a Next Best Offer (NBO) recommendation system, to be applied in the retail areas of banks.
    %\item Built predictive and segmentation models using high-dimensional data, with applications to financial decision-making. %Algorithms used: {\it Lasso, Random Forest, k-means, and xgboost}.
%\end{itemize}

%{\bf Ernst \& Young (EY) - Staff Data Scientist} \hfill {May 2023 - Jun 2024}
%\begin{itemize}
%\setlength\itemsep{-0.5em}
    \item Built quantitative workflows on high-dimensional administrative and financial data, including clustering, prediction, inference, and A/B testing for banking applications.
    %\item Design of segmentation models and subsequent classification of clients for a bank in Guatemala with more than 30 thousand clients. %Algorithms used: {\it LASSO, Random Forest, K-means, PCA, and classification trees}.
    %\item Designed an API using {\it FastAPI} and {\it SciPy} to infer loan interest rates from the financial system based on RCC data. The rate inferences achieved 94\% accuracy. These results are subsequently used for designing lead campaigns for financial institutions and banks.
   %\item Designed AB testing experiments in the largest retail bank in Guatemala.
    %\item Developed a library in {\it Rust} to increase the speed of pricing codes originally built in Python.
\end{itemize}

%{\bf BREIN - Data Science Trainee} \hfill {Sep 2022 - Apr 2023}
%
%\begin{itemize}
%\setlength\itemsep{-0.5em}
%    \item Participated in a catch prediction project for a fishing company using %{\it ARIMA} and {\it LSTM} models for time series. The model was designed to %process information automatically and produce monthly results.
%    \item Fuel optimization project for mining trucks using {\it xgboost} %algorithm. The constructed model was later integrated into Databricks.
%    %\item Correction of GPS variables using {\it Kalman filter}.
%\end{itemize}

{\bf Research Assistant to Prof. Noelia Bernal} \hfill {Sep 2021 - May 2022}
\begin{itemize}
\setlength\itemsep{-0.5em}
    %\item {\bf Options to expand contributory pension coverage among self-employed workers in Peru (\href{https://www.cepal.org/es/publicaciones/48051-%opciones-expandir-la-cobertura-contributiva-pensiones-%trabajadores}{ Bernal, 2022}):} Research consultancy on the analysis of the Peruvian pension system. Activities include: database management, descriptive statistical analysis, literature review, and report writing. 
    \item {\bf Does Matching Contribution incentivize informal workers to participate in retirement saving plans? A Randomized Evaluation in Peru (\href{https://doi.org/10.3386/w33925}{Bernal, Molina \& Galiani, 2025}):} Assistance in an impact evaluation project aimed at publication in an indexed journal. %Responsible for processing databases, creating graphs, regression analysis, and text processing in \LaTeX.
    %\item {\bf Can the pension coverage of self-employed workers in Latin America be improved? (\href{https://www.pinboxsolutions.com/resources/books/book-latin-america/}{Bernal \& Solf, 2023}):} Co-author of chapter in book 'Innovating Pensions for Self-Employed and Informal Workers in Latin America'.
\end{itemize} 
\end{rSection}

%----------------------------------------------------------------------------------------
%	TECHNICAL STRENGTHS SECTION
%----------------------------------------------------------------------------------------
\begin{rSection}{Awards and Recognitions}

{\bf Renzo Rossini Award for Research in Economics and Finance 2022-2023}\hfill {Sep 2023}
\\
First place at national level in the \href{https://www.bcrp.gob.pe/docs/Proyeccion-Institucional/premio-renzo-rossini/premio-renzo-rossini-2022-2023-resultados.pdf}{research contest} issued by the Central Reserve Bank of Peru for the research paper titled: "Monetary policy under a nominal GDP targeting rule: A small open economy approach".
\end{rSection}

\begin{rSection}{Publications and Presentations}

Can the pension coverage of self-employed workers in Latin America be improved? (\href{https://www.pinboxsolutions.com/resources/books/book-latin-america/}{Bernal \& Solf, 2023})
\\
Research Institute in Economics and Finance Annual Conference \href{https://ri-ef.com/rief-annual-conference-2023/}{2023}.
\\
Welfare implications of nominal GDP targeting in a small open economy (\href{https://mpra.ub.uni-muenchen.de/119999/}{Ortiz, Inca \& Solf, 2024})

\end{rSection}



\begin{rSection}{Relevant Skills}

\begin{tabular}{ @{} >{\bfseries}l @{\hspace{2ex}} l }
%Computing & Databricks, VBA, SQL (SQL Server), Azure, SAS.\\ 
Programming & Stata (proficient), Python (polars, scikit-learn, torch).  \\
Tools & Git, GitHub, SQL, Databricks, MLflow, FastAPI, LaTeX. \\
Languages &  Spanish (native), English (TOEFL iBT: 108 - C1), Italian (basic). 
\end{tabular}

\end{rSection}

\begin{rSection}{Teaching Experience}
{\bf Lecturer on MLOps program - Universidad Nacional de Ingeniería (Perú)} \hfill {Sep 2025 - Nov 2025}

%In charge of courses on MLOps. I taught methods to deploy models in production using FastAPI, Docker, and GCP. 
\end{rSection}



\begin{rSection}{Extracurricular Activities}
Fellow of the Economic Society for Bocconi Student. \\
%Fellow of the Quantitative Finance Society for Bocconi Students. \\
First place in DMC Hackathon - Aldeas Infantiles (2025). \\
%English teacher for children during summer classes. Program in collaboration with the Municipality of Lima.
%Participation in Hackathons organized by Rimac Seguros and Ferreyros (2024).

%I run a YouTube channel where I uploaded a \href{https://www.youtube.com/watch?v=k6m2h9tH7dE}{video about PCA}. I’m interested in gradually continuing this small project to explain Machine Learning concepts in Spanish through animations.



\end{rSection}




%----------------------------------------------------------------------------------------

%\begin{rSection}{References}


%\begin{tabular}{lr}
% Referee 1
%\begin{minipage}[t]{2.3in}
%\textbf{PhD Noelia Bernal}\\
%\textit{Assistant Professor} \\
%\textit{Centro de investigación de la Universidad del Pacífico}\\
%\textit{Universidad del Pacífico}\\
%\Telefon\ +(51) 941 455 593\\
%\Letter\ \href{mailto:n.bernall@up.edu.pe}{n.bernall@up.edu.pe}{\color{blue}}
%\end{minipage}

%\begin{minipage}[t]{2.5in}
%\textbf{PhD Oswaldo Molina}\\
%\textit{Investigador asociado} \\
%\textit{Executive Director} \\
%\textit{Centro de investigación de la Universidad del Pacífico}\\
%\textit{REDES}\\
%\Letter\ \href{mailto:o.molinac@up.edu.pe }{o.molinac@up.edu.pe }
%\end{minipage}

%\begin{minipage}[t]{2.5in}
%\textbf{PhD Marco Ortiz}\\
%%\textit{Investigador asociado} \\
%\textit{Assistant Professor} \\
%\textit{Universidad del Pacífico}\\
%\textit{REDES}\\
%\Letter\ \href{mailto:ma.ortizs@up.edu.pe}{ma.ortizs@up.edu.pe }
%\end{minipage}

%\begin{minipage}[t]{2.5in}
%\textbf{PhD Francisco Rosales}\\
%\textit{Investigador asociado} \\
%\textit{Data Science Chapter Lead} \\
%\textit{Centro de investigación de la Universidad del Pacífico}\\
%\textit{BREIN}\\
%\Letter\ \href{mailto:francisco@brein.pe}{francisco@brein.pe}
%\end{minipage}

%\end{tabular}

%\\

 
%\end{rSection}




\end{document}
